% Einheit 3 – Summe mal Summe (binomische-formeln.md Einheit 1) – Nr. 21–28
\einheitenkopf[e3][3 Summe mal Summe]{Einheit 3 von 5 · Summe mal Summe}
\zweigzeile{Hier lernst du, Terme mit zwei Variablen umzuformen und zwei Klammern miteinander malzunehmen · neu oder schon bekannt – je nach Buch · keine P10-Aufgabe · baut auf: Klammern auflösen (Einheit 1), Terme malnehmen}
\setcounter{aufgabe}{20}

\begin{aufgabe}{Ich erkenne, was mit was malgenommen wird.}
\anweisung{Zeichne Pfeile von jedem Glied der ersten Klammer zu jedem Glied der zweiten Klammer. Rechne nichts aus.}
\begin{teilezwei}
\tz \rule{0pt}{12mm}{\large $(x + 2) \cdot (y + 5)$} & \tz \rule{0pt}{12mm}{\large $(u + 4) \cdot (v - 1)$} \\
\tz \rule{0pt}{12mm}{\large $(x - 3) \cdot (x + 8)$} & \tz \rule{0pt}{12mm}{\large $(2x + 1) \cdot (x - 6)$} \\
\tz \rule{0pt}{12mm}{\large $(y + 7) \cdot (z + 3)$} & \\
\end{teilezwei}
\end{aufgabe}

\verfahren{Terme mit zwei Variablen}

\begin{aufgabe}{Ich kann den Wert eines Terms mit zwei Variablen berechnen.}
\anweisung{Setze ein und berechne den Wert des Terms.}
\begin{geruest}
\gz{a}{$2x + y$ \quad für $x = 3$ und $y = 4$}{\leerfeld}
\gz{b}{$x + 3y$ \quad für $x = 2$ und $y = 5$}{\leerfeld}
\gz{c}{$4x - y$ \quad für $x = 2$ und $y = 9$}{\leerfeld}
\gz{d}{$xy + 1$ \quad für $x = 3$ und $y = 4$}{\leerfeld}
\gz{e}{$3x - 2y$ \quad für $x = -2$ und $y = 5$}{\leerfeld}
\gz{f}{$x^2 + xy$ \quad für $x = -3$ und $y = 2$}{\leerfeld}
\gz{g}{$2x^2 - xy + y^2$ \quad für $x = -1$ und $y = 4$}{\leerfeld}
\end{geruest}
\end{aufgabe}

\begin{aufgabe}{Ich kann Terme mit zwei Variablen zusammenfassen (neu in diesem Jahr).}
\anweisung{Fasse zusammen.}
\begin{gleichungsraster}[1]
\gl{3x + 2y + 4x} & \gl{5u + v + 2v} \\
\gl{2x + 6y + 3x + 4y} & \gl{4u + 3v + 2u + 5v} \\
\gl{6x - y - 2x + 4y} & \gl{x^2 + 3xy + 2x^2 - xy + 5x} \\
\gl{4xy - 2x^2 + 3yx - x^2 + 6y} & \\
\end{gleichungsraster}
\end{aufgabe}

\begin{aufgabe}{Ich kann eine Klammer mit zwei Variablen auflösen.}
\anweisung{Löse die Klammer auf.}
\rechenplatz{2}
\begin{gleichungsraster}[1]
\gl{3 \cdot (x + 2y)} & \gl{4 \cdot (2x - y)} \\
\gl{5 \cdot (u + v)} & \gl{2 \cdot (3x + 4y)} \\
\gl{x \cdot (y + 3)} & \gl{-(2x - 5y)} \\
\gl{2x \cdot (x + 3y)} & \\
\anweisung{Löse die Klammern auf und fasse zusammen.}
\gl[2]{4 \cdot (x + 2y) - 3 \cdot (x - y)} & \\
\end{gleichungsraster}
\end{aufgabe}

\verfahren{Klammer mal Klammer}

\begin{aufgabe}{Ich kann zwei Klammern miteinander malnehmen.}
\anweisung{Multipliziere die Klammern aus.}
\rechenplatz{3}
\begin{gleichungsraster}[2]
\gl{(x + 2) \cdot (y + 4)} & \gl{(x + 5) \cdot (y + 1)} \\
\gl{(u + 3) \cdot (v + 6)} & \gl{(x + 7) \cdot (y + 2)} \\
\anweisung{Multipliziere aus und fasse zusammen.}
\gl{(x + 2) \cdot (x + 6)} & \gl{(x - 4) \cdot (x + 9)} \\
\gl{(x - 3) \cdot (x - 8)} & \gl{(2x + 1) \cdot (x + 7)} \\
\gl{(x + 2y) \cdot (3x + y)} & \gl{(3x - 2) \cdot (2x - 5)} \\
\end{gleichungsraster}
\end{aufgabe}

\begin{aufgabe}{Ich finde den Fehler beim Malnehmen zweier Klammern.}
\anweisung{Finde den Fehler, benenne ihn und korrigiere die Rechnung.}
\begin{teile}
\teil Paul rechnet: \rechnung{(x + 4) \cdot (x + 6) &= x^2 + 24}
\teil Lina rechnet: \rechnung{(x - 5) \cdot (x - 2) &= x^2 - 2x - 5x - 10 \\ &= x^2 - 7x - 10}
\teil Can rechnet: \rechnung{(x + 7) \cdot (x + 1) &= x^2 + x + 7x + 7 \\ &= x^2 + 15x}
\end{teile}
\end{aufgabe}

\begin{aufgabe}{Ich kann am Rechteck begründen, warum vier Produkte entstehen.}
Das Rechteck hat die Seitenlängen $x + 3$ und $x + 2$.

\flaechenbild{x}{3}{x}{2}{}{}{}{}

\begin{teile}
\teil Begründe am Bild, warum beim Malnehmen von $(x + 3) \cdot (x + 2)$ vier Produkte entstehen.
\teil Tim sagt: „$(x + 3) \cdot (x + 2) = x^2 + 6$.“ Begründe am Bild, was er vergessen hat.
\end{teile}
\end{aufgabe}

\begin{aufgabe}{Ich kann eine Sachaufgabe mit zwei Klammern lösen.}
Ein quadratisches Beet hat die Seitenlänge $x$\,m. Es wird neu angelegt: $3$\,m länger und $2$\,m schmaler.

\rechteck[seiten={x+3,x-2,,},punkte=]{4.5}{2.5}

\begin{teile}
\teil Stelle einen Term für den Flächeninhalt des neuen Beets auf und multipliziere aus.
\teil Um wie viel m$^2$ ändert sich der Flächeninhalt gegenüber dem alten Beet? Gib einen Term an.
\teil Wird ein Beet mit $x = 4$ dadurch größer oder kleiner? Und eines mit $x = 10$?
\end{teile}
\end{aufgabe}
